\chapter{REFERENCIAL TEÓRICO}
\label{chap:refteorico}

\section{Engenharia clínica e gestão de tecnologia médico-hospitalar}

A engenharia clínica é o campo responsável por aplicar conhecimentos de
engenharia ao ambiente de saúde, com foco na gestão segura e eficiente das
tecnologias médico-hospitalares ao longo de todo o seu ciclo de vida --- da
aquisição ao descarte. A gestão de tecnologia médico-hospitalar (GTMH) abrange
atividades como incorporação, inventário, manutenção, controle de desempenho e
desativação dos equipamentos, com o objetivo de assegurar disponibilidade,
confiabilidade e segurança. No Brasil, o gerenciamento de tecnologias em saúde
nos estabelecimentos é objeto de regulamentação sanitária específica.

\todo[inline]{Cite aqui a regulamentação e a literatura de engenharia clínica que
você adotar (ex.: ANVISA RDC nº 509/2021; livros/artigos de GTMH). Insira no
formato ABNT.}

\section{Manutenção e ciclo de vida dos equipamentos}

As estratégias de manutenção costumam ser classificadas em corretiva (executada
após a falha), preventiva (programada para reduzir a probabilidade de falha) e
preditiva (baseada no monitoramento da condição do equipamento). O histórico de
manutenção corretiva, em particular, é um indicador relevante do estado de um
equipamento: o aumento da frequência e, sobretudo, do custo das intervenções
tende a acompanhar o envelhecimento do item, sinalizando a aproximação do fim de
sua vida útil econômica. O acompanhamento desses indicadores ao longo do tempo
fornece subsídios para decidir entre manter ou substituir um equipamento.

\todo[inline]{Inclua referências sobre tipos de manutenção e ciclo de vida/curva
de custos de equipamentos.}

\section{Criticidade e priorização de substituição}

Nem todos os equipamentos têm o mesmo impacto sobre a assistência. A criticidade
expressa a importância clínica e o risco associado a cada equipamento, sendo
usualmente representada em níveis (por exemplo, de baixa a alta). A priorização
de substituição busca ordenar os candidatos à troca combinando múltiplos
critérios --- tipicamente criticidade, idade e custo de manutenção ---, de modo
que os recursos, sempre limitados, sejam direcionados aos itens cuja renovação
traz maior retorno em segurança e em eficiência. Abordagens multicritério desse
tipo são amplamente empregadas no apoio à decisão em gestão de ativos e
fundamentam o modelo de otimização adotado neste trabalho.

\todo[inline]{Cite a literatura de priorização/análise multicritério aplicada à
substituição de equipamentos.}

\section{Integração de dados (ETL) e visualização da informação}

O paradigma ETL (\textit{Extract, Transform, Load}) descreve o processo de
extrair dados de fontes diversas, transformá-los para um formato consistente e
carregá-los em um destino de análise. Trata-se de uma abordagem fundamental
quando os dados de origem são heterogêneos, como ocorre com planilhas de
diferentes sistemas e épocas. Uma vez integrados, os dados podem ser explorados
por meio de painéis interativos (\textit{dashboards}), recurso característico de
soluções de \textit{Business Intelligence} (BI), que sintetizam grandes volumes
de informação em indicadores e visualizações de fácil leitura, favorecendo a
tomada de decisão baseada em dados.

\todo[inline]{Cite referências de ETL, BI e visualização de dados que embasarão
esta subseção.}
